%% =====================================================================
%% appendix-proofs.tex -- proofs for formal.tex (Sections 4 and 5) and elision.tex (Section 8).
%% Defines: app:proofs, app:sound, app:abstraction, app:hostdom, app:guard, app:fast, app:precision, app:window,
%% app:elide, app:repair, app:counter.
%% =====================================================================
\providecommand{\Sc}[1]{\mathit{Sc}_{#1}}
\providecommand{\ASc}[2]{\mathit{ASc}_{#1}^{#2}}
\providecommand{\ESs}{\mathit{ES}^{\mathrm{src}}}
\providecommand{\ESd}{\mathit{ES}^{\mathrm{dst}}}
\providecommand{\AS}{\mathit{AS}}
\providecommand{\chainto}{\rightsquigarrow}
\providecommand{\mspecl}[1]{\ensuremath{\Rightarrow_{m}^{#1}}}
\providecommand{\ordl}[1]{\mathit{ord}_{#1}}
\providecommand{\Hpw}{\ensuremath{\Haz_C^{\mathit{pw}}}}
\providecommand{\Hep}{\ensuremath{\Haz_C^{\mathit{ep}}}}
\providecommand{\Hspec}{\ensuremath{\Haz_{\mathit{spec}}}}
\providecommand{\pc}[1]{\textsc{#1}}
\providecommand{\reach}{\hookrightarrow}
\providecommand{\posof}{\pi}
\providecommand{\Hstrict}{\ensuremath{\Haz_{\mathit{spec}}^{\mathit{dom}}}}

\section{Proofs}
\label{app:proofs}

The proofs refer to the algorithm of Figure~\ref{fig:alg} by procedure and line. Throughout, $\tau$ is fixed and the
algorithm runs in the verified configuration (no relaxation for transitions from \code{UNDEFINED}). An event is
\emph{processed} when \pc{Access} is called for it.

\paragraph{Notation}
For an event $e$ and a link $(D,c)$ we write $e \reach_x (D,c)$ if there is a chain $(D_1,c_1) \chainto \dots
\chainto (D_k,c_k) = (D,c)$ ($k \ge 1$) with $e \in \Sc{1}(D_1)$ and $c_1$ given by the class rule of
Definition~\ref{def:spec}, and $e \reach_m (D,c)$ if moreover $e \in \ASc{1}{\ell}(D_1)$ for every $\ell \in R(e)$. We
write $e \reach (D,c)$ for $e \reach_m (D,c)$ when $k(e) = \mathsf{W}$ and for $e \reach_x (D,c)$ otherwise, and drop
the class when it is $\mathsf{G}$. If $e \reach (D,c)$ and $(D,c) \chainto (D',c')$ then $e \reach (D',c')$, because
the intermediate dependencies of a chain need only chain. For a signal or a wait, the access scopes equal the
synchronization scopes, so $e \in \Sc{1}(\mathsf{Sig})$ gives $e \reach_m \mathsf{Sig}$. We say that a link $(D,c)$
\emph{holds} $u$ for source $e$ if $s(u)$ is in $D$'s second stage set restricted by $c$, and $\posof(u) = \posof(e)
\ne \bot$ when $c = \mathsf{L}$. Two facts about classes are used throughout. (C1)~A $\mathsf{G}$ link dominates the $\mathsf{N}$ and
$\mathsf{L}$ links of the same dependency: its restricted second stage set contains theirs and it holds without a
token condition, so it holds every access they hold, and every link they chain into is chained into from it in class
$\mathsf{G}$, which again dominates. (C2)~Signals and waits have class $\mathsf{G}$, whatever class the link before
them has.

Two facts about the trace translation are used repeatedly. (F1)~The structs of one barrier call are consecutive on
their queue, and only their transitions lie at their positions (W2). (F2)~A semaphore wait has an empty first stage set
and a signal an empty second stage set, so a barrier never chains into a wait, and a chain leaves a queue only through
a signal and enters one only through the paired wait.

%% -------------------------------------------------------------------
\subsection{Soundness (Theorem~\ref{thm:sound})}
\label{app:sound}

\begin{lemma}[Linearization]\label{lem:lin}
Positions on a queue are assigned in processing order, a wait is processed after its signal, and host operations are
processed in host order. Hence if $D \chainto D'$ then $D$ is processed before $D'$, and if $a \xspec b$ then $a$ is
processed before $b$.
\end{lemma}
\begin{proof}
Line~7 of \pc{Check} assigns $p(o) = \mathit{VC}[n][n]$ after incrementing it, once per operation, in queue order.
Line~3 processes a wait only after its name has been signalled. A chain step is either a step on one queue to a later
position, a step on the host to a later position, or a signal to its paired wait (W1 makes the pairing unique).
Each goes forward in processing order, and so does the first step $a \in \Sc{1}(D_1)$ ($p(a) < p(D_1)$) and the last
$b \in \Sc{2}(D_n)$ ($p(D_n) < p(b)$).
\end{proof}

\begin{lemma}[Chain state]\label{lem:chain}
Let \pc{Chain}$(e)$ have been advanced over the barriers of $q(e)$ up to some point. Call a link $(B,\mathsf{G})$
\emph{x-reached} if $g_x$ held at $B$ (line~7 of \pc{Advance}), $(B,\mathsf{N})$ x-reached if $l_x$ held, and
$(B,\mathsf{L})$ x-reached if $l_x$ held and $\lambda(B) = \mathsf{loc}$ (line~8); m-reached likewise with $g_m$,
$l_m$ (line~9).
(a)~If $(B,c)$ is x-reached then $e \reach_x (B,c)$ through a chain of barrier structs on $q(e)$.
(b)~If $(B,c)$ is m-reached then $e \reach_m (B,c)$, and $e$ is in the first access scope of the first struct of that
chain on every element of $R(e)$.
(c)~Conversely, if $e \reach_x (B,c)$ through barrier structs of $q(e)$ only, then $(B,c)$ or $(B,\mathsf{G})$ is
x-reached, and the same for m-reached.
For the chain of a wait record, the same statements hold with ``$e \in \Sc{1}(D_1)$'' replaced by ``$W \chainto D_1$''
for the wait $W$ (all its links have class $\mathsf{G}$), and the m-reached and x-reached sets coincide.
\end{lemma}
\begin{proof}
Invariant: when \pc{Advance} reaches struct $B$, $X_o$ is the union, over the x-reached links of earlier calls, of
their second stage sets restricted by class ($\mathsf{G}$ and $\mathsf{N}$), $X^L_o$ the same for the $\mathsf{L}$
links, $X_c[s]$ and $X^L_c[s]$ those over the current call with struct identity $s$, and $\mathit{snap}[v]$ the four
unions at the latest set of event $v$ since the chain started, deleted by a reset. The same holds for the $M$ sets and
m-reached links, whose ends $M_{\mathit{ends}}$ ($\mathsf{G}$, and $\mathsf{N}$ restricted to stages outside $F$) and
$M^L_{\mathit{ends}}$ ($\mathsf{L}$) keep the second access masks. The invariant holds initially (all sets empty, or
$X_o = M_o = E$ for a wait record, which stands for the wait's second stage set in class $\mathsf{G}$) and is kept by
lines~1--3 and~6--9, since a new call folds the current one into the earlier ones. Advancing incrementally is a left
fold over the barrier list, so it computes the same sets as one pass.

(a) by induction over the order in which links are reached. Write $X_f$, $X^L_f$ for the sets of line~3.
\emph{Base:} $d_x$ holds, so $e \in \Sc{1}(B)$: line~4 tests $p(e) < p(B)$ (the chain starts at the first barrier at
or after $p(e)$, and a transition is admitted only at its own struct), the stage, the call condition and, for an
event wait, that $\mathit{snap}[v]$ exists, which holds exactly when a set of $v$ lies between $e$ and $B$ without a
later reset, that is $p(e) < c_B$. The class rule of Definition~\ref{def:spec} enters $B$ locally exactly when
$\lambda(B) \ne \mathsf{glob}$ and $e$ is located, which is $\mathit{lo}$; line~6 then sets $l_x$ (classes
$\mathsf{N}$, and $\mathsf{L}$ if $\mathsf{loc}$) and otherwise $g_x$ (class $\mathsf{G}$). \emph{Step:} $X_f$ or
$X^L_f$ meets $\ESs(S_1)$. Then some x-reached link $(B',c')$ has its restricted second set meeting $\ESs(S_1)$, with
$p(B') < p(B)$, $B'$ not another struct of $B$'s call, and $p(B') < c_B$ for an event wait, so $(B',c') \chainto
(B,c)$ for the class $c$ that the rule gives: from $X^L_f$ ($c' = \mathsf{L}$) into a local $B$, classes $\mathsf{N}$
and $\mathsf{L}$, which is $l_x$; from $X_f$ ($c' \in \{\mathsf{G},\mathsf{N}\}$) or into a global $B$, class
$\mathsf{G}$, which is $g_x$. By induction $e \reach_x (B',c')$.

(b) is the same argument for line~9. Its base case $d_m$ adds to $d_x$ that $k(e) = \mathsf{W}$, $s(e) \in
\AS(S_1)$, $f(e) \in_{\mathsf{W}} A_1$ and $r_B = \ast$ or ($r_B = r(e)$ and $R_B \supseteq R(e)$), which is $e \in
\ASc{1}{\ell}(B)$ for every $\ell \in R(e)$. A transition is in both access scopes of its own barrier.

(c) by induction on the length of the chain, with the same case split read backwards: if $(B',c') \chainto (B,c)$ and
$(B',c')$ is x-reached, its restricted second set is in $X_f$ or $X^L_f$ (by the invariant; same call and struct, or
an earlier call, or the snapshot of an event wait with $p(B') < c_B$), and it meets $\ESs(S_1)$, so line~6 fires for
the class the rule gives. If only $(B',\mathsf{G})$ is x-reached, (C1) gives $(B,\mathsf{G})$.

For a wait record, $d_x = d_m = \mathit{false}$, the seeds of $X$ and $M$ are equal and the local sets start empty,
so only class $\mathsf{G}$ arises and lines~7 and~9 fire together.
\end{proof}

\begin{lemma}[Same queue]\label{lem:same}
Let $q(e) = q(u)$ and let line~4 or line~5 of \pc{Ordered} return \emph{ok}. Then $\ordl{\ell}(e,u)$ holds for every
$\ell \in O = R(e) \cap R(u)$.
\end{lemma}
\begin{proof}
All barriers in $\mathit{Bars}[n]$ precede $u$. \emph{Line~5}, execution mode: $s(u) \in X_{\mathit{tot}}$, so some
x-reached link $(B,c)$ with $c \in \{\mathsf{G},\mathsf{N}\}$ has $s(u)$ in its restricted second set (for $s(u) =
\bot$, only a $\mathsf{G}$ set of $\top$ contains it), hence it holds $u$ and $e \reach_x (B,c)$ by
Lemma~\ref{lem:chain}(a); or $\posof(u) = \posof(e) \ne \bot$ and $s(u) \in X^L_{\mathit{tot}}$, so an x-reached
$(B,\mathsf{L})$ holds $u$. Memory mode: some m-reached link holds $u$ with $u \in \ASc{2}{\ell}(B)$ for every $\ell
\in O$ (the test of \pc{Reaches}: $r_B = \ast$ or $R_B \supseteq O$, stage in the restricted set and in $\AS(S_2)$,
flag in $A_2$, and the token for an $\mathsf{L}$ end), and $e \reach_m (B,c)$ with $e$ in the first access scope on
all of $R(e) \supseteq O$ (Lemma~\ref{lem:chain}(b)).
\emph{Line~4}, $u = t_B$: $u$ is processed when $B$ is appended, so $B$ is the last barrier. If the chain reached $B$
last globally, then $e \reach (B,\mathsf{G})$ or $e \reach (B,\mathsf{N})$ and $t_B$ is in both scopes of its own $B$;
a transition is located, so an $\mathsf{N}$ link alone does not hold it, but the chain reached $B$ in class
$\mathsf{G}$ there, since line~8 adds no stage of $\mathsf{LT}$ to the global sets and the last global position is set
only by line~7. If it reached $B$ last locally, then $\posof(e) = \posof(u)$ and the $\mathsf{L}$ link holds $t_B$.
Otherwise $e \in \Sc{1}(B)$ directly (the test repeats the conditions of $\Sc{1}$, including the event cut and the
exclusion of other structs of the call and of transitions), in memory mode $e \in \ASc{1}{\ell}(B)$ on $O$, and, when
$B$ is entered locally from $e$, $\lambda(B) = \mathsf{loc}$ and $\posof(e) = \posof(u)$, so the one-link chain
$(B,\mathsf{L})$ holds $t_B$; when $B$ is not entered locally the link is $(B,\mathsf{G})$.
\end{proof}

\begin{lemma}[Relays]\label{lem:clock}
Assume W1. Let $\mathsf{Sig}(x,S)$ be processed on queue $q_s$ with relay $\rho$, let $e$ be an event on queue $m$,
and let $(t, \varphi) \in \rho[m]$ with $p(e) \le t$, where $\varphi = \bot$ or \pc{InSignal}$(e,\varphi)$ holds. Then
$e \reach \mathsf{Sig}$.
\end{lemma}
\begin{proof}
By induction on processing order, over the signals. An entry of $\rho$ is never created by $\sqcup$, which only keeps
or drops entries, so every entry of $\rho[m]$ was put there by one of the lines of \pc{Relay}$(q_s, p(\mathsf{Sig}),
S)$. \emph{Line~1}, the signal's own entry ($m = q_s$, $t = p(\mathsf{Sig})$): $e$ is an operation of $q_s$ other than
$\mathsf{Sig}$, so $p(e) < p(\mathsf{Sig})$. If $S = \top$, $e \in \Sc{1}(\mathsf{Sig})$ (every stage, $\bot$ and
transitions included), which is also the signal's first access scope. Otherwise $\varphi$ describes $\mathsf{Sig}$ and
Lemma~\ref{lem:insig} below gives $e \reach \mathsf{Sig}$. \emph{Lines~2--4}, an entry of the relay $\rho'$ of a wait
$W$ on $q_s$ processed before $\mathsf{Sig}$, whose signal $\mathsf{Sig}'$ was processed earlier still. By induction,
$e \reach \mathsf{Sig}'$, and $\mathsf{Sig}' \chainto W$ by W1. It remains to show that $W$ chains into $\mathsf{Sig}$,
possibly through barriers of $q_s$. Line~2: $S = \top$, and $W$'s second stage set has a stage other than $\mathsf{LT}$
(line~5 of \pc{Check} keeps no other wait), so it meets $\top$ and $W \chainto \mathsf{Sig}$. Line~3: $W$ has $E =
\top$, and $\ESs(S)$ has a stage other than $\mathsf{LT}$, so they meet. Line~4: $\pc{Chain}(W).X_{\mathit{tot}}$ meets
$\ESs(S)$. Either the seed $E$ meets it, and $W \chainto \mathsf{Sig}$; or some link $(B,\mathsf{G})$ that the wait
record's chain x-reaches has $\ESd(S_2)$ meeting $\ESs(S)$, so $W \reach_x B$ by Lemma~\ref{lem:chain} (wait-record
form) and $B \chainto \mathsf{Sig}$. The intermediate dependencies of a chain need only chain (Definition~\ref{def:spec}), so in
every case $e \reach \mathsf{Sig}$, and in memory mode the first link is still the one that holds $e$ in its first
access scope.
\end{proof}

\begin{lemma}[Signal filter]\label{lem:insig}
If \pc{InSignal}$(e,\varphi)$ holds for the masked signal $\mathsf{Sig}$ that $\varphi$ describes, and $e$ precedes it
on its queue, then $e \reach \mathsf{Sig}$.
\end{lemma}
\begin{proof}
Line~1: $e$ is not a transition and $s(e)$ is in the signal's first stage set, so $e \in \Sc{1}(\mathsf{Sig})$, which is
also its first access scope. Line~3: the chain of $e$ is advanced over the barriers before the signal. For a write,
some m-reached link $(B,c)$, of any class, has its restricted second stage set meeting the signal's first stage set,
so $(B,c) \chainto \mathsf{Sig}$, a link of class $\mathsf{G}$ by (C2), and $e \reach_m (B,c)$
(Lemma~\ref{lem:chain}(b)). For a read, the same with an x-reached link. This is why \pc{InSignal} tests the union of
the global and the local sets: a local chain followed by a signal is an ordinary chain.
\end{proof}

\begin{lemma}[Waits]\label{lem:viawait}
Assume W1. If \pc{ViaWait} returns true for $(e,u,\mathit{mem},O)$, then $\ordl{\ell}(e,u)$ holds for every $\ell \in
O$.
\end{lemma}
\begin{proof}
Let $m = q(e)$, $n = q(u)$. \emph{Line~1.} $\mathit{Full}[n][m] \ge p(e)$ means that some wait $W$ on $n$ processed
before $u$ has $E = \top$ and an unfiltered entry $(t, \bot) \in W.\rho[m]$ with $t \ge p(e)$ (line~6 of \pc{Check}).
By Lemma~\ref{lem:clock}, $e \reach \mathsf{Sig} \chainto W$, and $u \in \Sc{2}(W)$ because $W$ has every stage. A
transition is in $\Sc{2}$ of an earlier wait of every stage, and the access scope of a wait is its synchronization
scope. No other shortcut exists: in particular, a wait on a semaphore signalled by the host is accepted here only if
it has $E = \top$. \emph{Lines~2--6.} The record $w$ of a wait $W$ is used only if (line~3) its relay holds an entry
$(t,\varphi)$ for $m$ with $t \ge p(e)$ and $\varphi = \bot$ or \pc{InSignal}$(e,\varphi)$. Lemma~\ref{lem:clock}
then gives $e \reach \mathsf{Sig} \chainto W$, so $e \reach W$. Then: line~4, $u \in \Sc{2}(W)$ as above. Line~5, $u =
t_B$ and $W$'s chain reaches $B$, the last barrier on $n$, so $W \reach B$ (Lemma~\ref{lem:chain}) and $t_B$ is in both
scopes of $B$. Line~6, $s(u) \in E$ gives $u \in \Sc{2}(W)$, or $W$'s chain reaches $u$ as in Lemma~\ref{lem:same},
line~5. In memory mode the first link of the chain from $e$ is a signal, or a barrier that covers $R(e)$
(Lemma~\ref{lem:insig}), and the last link holds $u$ in its second access scope on all of $O$.
\end{proof}

\begin{proof}[Proof of Theorem~\ref{thm:sound}]
Let \pc{Ordered}$(e,u)$ return \emph{ok}, with $\mathit{mem} = (k(e) = \mathsf{W})$, which is the choice of
$\ordl{\ell}$. Line~2: both are on the host, and host order is an execution and memory order (A-HOST). Lines~4--5:
Lemma~\ref{lem:same}. Lines~6 and~8: Lemma~\ref{lem:viawait}. So $\ordl{\ell}(e,u)$ holds for all $\ell \in O$, that
is, $e \spec u$.

For the inclusion, let $\{a,b\}$ race at $\ell$, and let $e$ be the one that is processed first and $u$ the other. In pairwise
mode $e$ is never removed from any segment, and $u$'s segments partition $R(u)$, so $e$ is in the segment that holds
$\ell$. It is a candidate of \pc{Access}$(u)$: writes always are, and reads are when $u$ writes, which a conflict
requires. The pair is not exempt by E1--E4, so \pc{Ordered}$(e,u)$ is called. It cannot return \emph{ok}, since that
would give $\ordl{\ell}(e,u)$. Line~4 of \pc{Access} does not apply, because an E5 pair is exempt
and does not race. So line~5 reports the pair.

The locality of Definition~\ref{def:spec} enters this proof only through Lemmas~\ref{lem:chain} and~\ref{lem:same}: the
chain state distinguishes the three classes exactly as the definition does, and a local ending is accepted only with
equal tokens. Lemma~\ref{lem:clock} needs no change beyond (C2), because every relay entry passes through a signal, and
Lemma~\ref{lem:viawait} none, because a wait record's chain starts in class $\mathsf{G}$.
\end{proof}

%% -------------------------------------------------------------------
\subsection{From the run to the events (Corollary~\ref{cor:sched})}
\label{app:abstraction}

Model the run as the trace $\tau$ (A1) with its \emph{true} accesses $\alpha$ in place of the events. Each $\alpha$
has its exact stage, flag and byte or element set. The dependencies are those of $\tau$, and the true accesses are
ordered by Definition~\ref{def:spec} applied to them. A device access $\alpha$ is \emph{represented} by an event $a$ (A2)
if they have the same queue and position (that of the command), $R(a) \supseteq R(\alpha)$, $s(a) \in \{s(\alpha),
\bot\}$, $f(a) \in \{f(\alpha), \bot\}$, $k(a) = \mathsf{W}$ if $\alpha$ writes, and every exemption tag of $a$ is
also a true exemption of $\alpha$. A CPU store $\sigma$ is represented by a host write $w$ as in A2-H.

\begin{lemma}[Abstraction]\label{lem:abs}
If $a$ represents $\alpha$ and $b$ represents $\beta$, then for $\ell \in R(\alpha) \cap R(\beta)$,
$\ordl{\ell}(a,b)$ implies $\ordl{\ell}(\alpha,\beta)$. Hence, if $\alpha$ and $\beta$ race at $\ell$, then $a$ and
$b$ race at $\ell$.
\end{lemma}
\begin{proof}
Scope membership is monotone in the direction we need. $\bot$ lies only in a set of every stage, which contains
$s(\alpha)$. A $\bot$ flag lies only in a mask that contains \code{MEMORY\_READ} or \code{MEMORY\_WRITE}, which contains
$f(\alpha)$. Access-scope membership at an element $\ell$ depends on $\ell \in R_B$, not on the rest of the region, and
positions are equal. So $a \in \Sc{i}(D)$ implies $\alpha \in \Sc{i}(D)$, and $a \in \ASc{i}{\ell}(D)$ implies
$\alpha \in \ASc{i}{\ell}(D)$, and a chain for $a,b$ is a chain for $\alpha,\beta$. If $\alpha$ writes, so does $a$, so
$\ordl{\ell}(a,b)$ uses the memory relation whenever $\ordl{\ell}(\alpha,\beta)$ does, and the memory relation implies
the execution relation. For a store $\sigma$ represented by $w$, A2-H states directly that $w$ and $\sigma$ have the
same predecessors and successors. It holds for page capture because $w$ sits between the two host calls that
bracket $\sigma$, where the host queue has no other operation, and host-queue orderings depend only on the position
relative to host signals and waits. For the race: $\ell \in R(a) \cap R(b)$, one of $a,b$ writes, and they are not
exempt, because a claimed exemption would be a true one. Neither $\ordl{\ell}(a,b)$ nor $\ordl{\ell}(b,a)$ holds,
since either would order $\alpha,\beta$.
\end{proof}

\begin{proof}[Proof of Corollary~\ref{cor:sched}]
By Lemma~\ref{lem:abs}, a race of covered accesses gives a race in $\Hspec(\tau)$, which is empty. The pairwise case
follows from Theorem~\ref{thm:sound} and the epoch case from Lemma~\ref{lem:fast}. For a device schedule consistent
with $\tau$'s synchronization, an ordered pair executes in that order with the required availability and visibility,
by the specification's definition of the dependencies.
\end{proof}

%% -------------------------------------------------------------------
\subsection{Index guards (Lemma~\ref{lem:guard})}
\label{app:guard}

The specification counts a descriptor as accessed only if an invocation \emph{dynamically uses} it, that is, evaluates
an access chain into the array with its index; the layout and synchronization requirements attach to that access. So
it suffices to show that in every schedule consistent with $\tau$, every invocation of $d$ indexes the array only with
elements of $I(d)$.

\begin{proof}[Proof of Lemma~\ref{lem:guard}]
Let $V$ be the attribute bytes that $d$ fetches (the words of vertices or instances $\mathit{first}, \dots,
\mathit{first}+n-1$ at the recorded offset and stride), let $r$ be $d$'s vertex-input read, which covers $V$ and
carries $g$, and let $S$ be the position of $d$'s submission in host order.

\emph{Step 1: in every schedule the fetch returns the bytes rebuilt at $S$.} By C5 no device command writes $V$ and no
image is bound on it, so the only writes to $V$ are host stores, which A2-H represents as host write events $w$ at the
call boundary of each store, over every byte of each written page. Consider such a $w$ on $V$. If $w$ and $r$ are
unordered, they race, so the race is reported: in pairwise mode by Theorem~\ref{thm:sound}; in epoch mode against $r$,
or against a later read that stands in for $r$ and inherits its guard, or against a host write that $r$ was compared
with and that $w$ precedes in host order, by Lemma~\ref{lem:fast}. A reported pair with $r$ or its stand-in fails $g$
(C6), contradicting the premise; a reported pair of $r$ and an earlier host write likewise. If $w$ is ordered before
$r$ and recorded before $S$, its bytes are in the rebuilt contents, since contents are rebuilt from every record before
$S$ in host order (one host thread, P3, checked per memory by C4). If $w$ is ordered before $r$ and recorded after
$S$ (for example through a host-signalled timeline value), then when $r$ is checked, $w$ or a later host write that
supersedes it is a live write entry on $V$ recorded after $S$, which fails $g$ (C6). If $r$ is ordered before $w$, $w$
does not affect what $r$ reads. Hence the fetch returns the rebuilt bytes, which C4 requires to be recorded for all of
$V$.

\emph{Step 2: the shader maps them into $I(d)$.} By C1 and C2 the index of every access chain into the array is, on
every path, a constant or the index chain applied to the attribute value delivered to the invocation: a flat fragment
input receives the provoking vertex's output, which the vertex shader sets to the chain applied to the attribute on
every path (an output not set on some path, or set through a derived pointer, makes the index unknown, and the binding
is widened). The analysis follows local variables and pointer parameters through stores, loads and calls and joins
over all paths, ignoring branch conditions, which can only add values. Attribute fetch of a 32-bit integer format at
the recorded stride, offset and input rate of a non-indexed direct draw (C2, C3) is a function of $V$. So every index
lies in $I(d)$, and by C7 none is out of range.

\emph{Step 3: the dropped events stand for no access.} The elements outside $I(d)$ are not dynamically used in any
schedule, so their events correspond to no access of the run, and the remaining events still represent every access
as A2 requires. Lemma~\ref{lem:abs} and Corollary~\ref{cor:sched} apply to the narrowed event set.

\emph{No circularity.} Whether $g$ holds depends on the checker's verdicts, which depend on the event set, which
depends on which guards hold. The conditions C5 and C6 mention only events on $V$, which are buffer bytes, and the
happens-before order. The events a guard adds when it fails are image reads (C8), and accesses create no dependencies
in Definition~\ref{def:spec}. So neither the events on $V$ nor any ordering changes when conditional reads are added or
removed, and the set of failing guards is the same in every run; the checker reaches it by re-running with the failed
guards widened. The executable reference defines $\Hspec$ with the same rule, and the property tests of
$\Hspec \subseteq \Hpw$ include guarded frames.
\end{proof}

%% -------------------------------------------------------------------
\subsection{Host-domain operations (Proposition~\ref{prop:hostdom})}
\label{app:hostdom}

\begin{proof}[Proof of Proposition~\ref{prop:hostdom}]
Since no error is reported, Theorem~\ref{thm:sound} and Lemma~\ref{lem:fast} give $\Hspec(\tau) = \emptyset$. Let $e$ be a
device write and $u$ a later host write that conflict; then $e \spec u$, and we show that $e$ is host-domain ordered
before $u$.

\emph{Pairwise mode.} \pc{Access}$(u)$ compares $u$ with $e$, \pc{Ordered} returns \emph{ok}, and since no warning is
reported \pc{HostDomain}$(e,u)$ holds: the fresh chain of $e$, advanced over the barriers of $q(e)$ up to $t =
\mathit{Full}[\mathit{host}][q(e)]$, m-reaches a link $(B,c)$ that holds a host access at \code{HOST\_READ} or
\code{HOST\_WRITE} on all of $O$. By Lemma~\ref{lem:chain}(b), $e \mspecl{\ell} (B,c)$ for every $\ell \in O$, and $B$
has $\mathsf{HOST}$ in its second stage and access stage sets and a host access in its second access mask ($\mathsf{HOST}$
is outside $F$, so the class does not matter). $\mathit{Full}[\mathit{host}][q(e)] = t$ means that a host wait $W$
before $u$ has an unfiltered relay entry $(t,\bot)$ for $q(e)$, so by Lemma~\ref{lem:clock} every operation of $q(e)$
at a position up to $t$, $B$ among them, is in the first scope of a chain of signals and waits of every stage that ends
at $W$; and $u \in \Sc{2}(W)$, because a host wait holds every later host event. This is the chain from $B$ that
Definition~\ref{def:hostdom} asks for.

\emph{Epoch mode.} Suppose $e$ is not compared with $u$. Then $e$ was dropped before $u$ was processed, by an access $v$
that stands in for it, and we show by induction on drops that $v$ carries the domain operation over to $e$. If $v$ is
a host write, it dropped $e$ only with \pc{Ordered}$(e,v)$ = \emph{ok} and no warning (a warning marks the pair
\emph{premise}, which keeps $e$), so $e$ is host-domain ordered before $v$ by the pairwise argument, and the chain
from $B$ to $v$'s host wait also reaches $u$, which follows $v$ in host order. If $v$ is a device write ordered after
$e$ (D1), or a write that supersedes $e$ (D2), then, by the induction, $v$ is host-domain ordered before $u$ through
some $B$, possibly on another queue. For D1, $e \mspecl{\ell} v$ and $v \mspecl{\ell} (B,c)$ compose as in
Lemma~\ref{lem:trans}, with $B$ treated as an access at stage $\mathsf{HOST}$: $e \mspecl{\ell} (B,c)$. For D2,
Lemma~\ref{lem:subst} gives $e \mspecl{\ell} (B,c)$ directly. A transition $v$ is a D1 drop. Shadow entries are
compared with $u$ as in Lemma~\ref{lem:fast}. In every case $e$ is host-domain ordered before $u$.
\end{proof}

The test is conservative in two ways that the property test of the checker exercises: a domain operation on another
queue is credited only through a stand-in, and in pairwise mode not at all; and a domain operation reached only
through stage-masked signals is not credited. On random frames with one and two device queues, every pair that the
pairwise mode warns about has its host write flagged (by a warning or a hazard) in epoch mode, and the epoch mode adds
no warning.

%% -------------------------------------------------------------------
\subsection{The epoch mode (Lemmas~\ref{lem:modes} and~\ref{lem:fast})}
\label{app:fast}

\begin{proof}[Proof of Lemma~\ref{lem:modes}]
\pc{Ordered} reads only the clocks, wait records, barrier lists and chains, and none of them depends on the
mode. Every candidate of the epoch mode, including shadow and $\mathit{hid}$ entries, is an earlier event that
overlaps $u$ and conflicts with it, so it is also a candidate in pairwise mode. The exemption test and the R-LT
test are the same in both modes.
\end{proof}

The epoch mode removes an entry $e$ from a segment $\varsigma \subseteq R(u)$ in three ways. (D1)~The write $u$ is
ordered after it: \pc{Ordered}$(e,u) = $ \emph{ok}. (D2)~The write $u$ supersedes the write $e$. (D3)~The read $u$
replaces the read $e$ of the same key $\kappa$. The proof needs two facts about the definition.

\begin{lemma}[Transitivity]\label{lem:trans}
Let $\ell \in R(a) \cap R(v) \cap R(x)$. If $a \xspec v$ and $v \xspec x$ then $a \xspec x$. If $a \mspecl{\ell} v$ and $v
\mspecl{\ell} x$ then $a \mspecl{\ell} x$.
\end{lemma}
\begin{proof}
Host order is transitive, and a host event before (after) the first (last) host dependency of a chain is in its
scope, so host-order steps compose with chains. Otherwise take chains $C_1 = D_1 \dots D_n$ for $(a,v)$ and $C_2 =
D'_1 \dots D'_k$ for $(v,x)$, with $v \in \Sc{2}(D_n) \cap \Sc{1}(D'_1)$. In memory mode $a$ is in the first access
scope of $D_1$ and $x$ is in the second access scope of $D'_k$ at $\ell$.

\emph{$v$ is not a transition.} $D_n$ and $D'_1$ are on $q(v)$ with $p(D_n) < p(v) < p(D'_1)$. Both stage sets contain
$s(v)$, or both are $\top$ if $s(v) = \bot$, so they meet. By (F1), no access lies between two structs of one call, so
they are not two structs of one call. If $D'_1$ is an event wait, $p(D_n) < p(v) < c_{D'_1}$. Hence $D_n \chainto
D'_1$ and $C_1C_2$ is a chain.

\emph{$v = t_B$.} By Definition~\ref{def:spec}, $D_n$ is $B$ or a wait $W$ of every stage before $B$, and $D'_1$ is $B$
or a signal $\mathsf{Sig}$ of every stage after $B$. (i)~$D_n = D'_1 = B$: concatenate and share $B$. (ii)~$D_n = W$ and
$D'_1 = \mathsf{Sig}$: $W$'s second set is $\top$ and meets $\top$, so $W \chainto \mathsf{Sig}$. (iii)~$D_n = W$ and
$D'_1 = B$: if $k \ge 2$ then $D'_2$ follows $B$ on $q(B)$ (by F2, $D'_2$ is not a wait) or is a signal there. Its
first stage set meets $B$'s second set, so it has a stage other than $\mathsf{LT}$ and $W \chainto D'_2$; event cuts
carry over, since $p(W) < p(B)$. Use $C_1 D'_2 \dots D'_k$. If $k = 1$, then $x$ follows $B$ on $q(B)$, so $x \in
\Sc{2}(W)$, which equals $\ASc{2}{\ell}(W)$, and $C_1$ alone orders $a$ before $x$. (iv)~$D_n = B$ and $D'_1 =
\mathsf{Sig}$: if $n \ge 2$, $D_{n-1} \chainto B$, so $D_{n-1}$'s second set has a stage other than $\mathsf{LT}$ and
$D_{n-1} \chainto \mathsf{Sig}$ ($D_{n-1}$ precedes $B$ on $q(B)$ by F2), so use $D_1 \dots D_{n-1} C_2$. If $n = 1$,
$a$ precedes $B$ on $q(B)$, so $a \in \Sc{1}(\mathsf{Sig})$, which equals the signal's first access scope, and $C_2$
alone orders $a$ before $x$. In every case the first link keeps $a$ in its first (access) scope and the last link keeps
$x$ in its second, so the memory statement follows too.

\emph{Classes.} It remains to check that the concatenated chain has classes that still hold $x$ for source $a$. The
chain for $(v,x)$ starts at $D'_1$ in the class that $v$ enters it with: locally if $D'_1$ is local and $v$ is
located. The chain for $(a,v)$ ends at $D_n$ in a class $c_n$ that holds $v$. If $c_n = \mathsf{G}$, the
concatenation enters $D'_1$ in class $\mathsf{G}$, which by (C1) dominates whatever class $v$ entered it with, so
every later link of $C_2$ has a dominating counterpart and $x$ is held without a token condition. If $c_n =
\mathsf{N}$, $s(v) \notin F$, so $v$ is not located and enters $D'_1$ in class $\mathsf{G}$, as does the
concatenation. If $c_n = \mathsf{L}$, then $\posof(v) = \posof(a) \ne \bot$ and $v$ is located; the concatenation
enters $D'_1$ from an $\mathsf{L}$ link, which by the class rule gives the same classes as entering it from the located
event $v$, so $C_2$'s classes are unchanged, and an $\mathsf{L}$ ending of $C_2$ holds $x$ for source $v$, that is with
$\posof(x) = \posof(v) = \posof(a)$, so it holds $x$ for source $a$. In the transition cases (i)--(iv) the shared or
bridging dependency is $B$ itself, a wait or a signal; waits and signals are $\mathsf{G}$ by (C2), and for (i) the
same argument applies at $B$.
\end{proof}

\begin{lemma}[Substitution]\label{lem:subst}
Let $e$ and $u$ be events on one queue with $p(e) < p(u)$, equal stage, kind and position token, neither a transition
nor in a barrier call, and let $\ell \in R(e) \cap R(u)$. If $u \xspec x$ then $e \xspec x$. If moreover $f(e) = f(u)$ and $u
\mspecl{\ell} x$, then $e \mspecl{\ell} x$.
\end{lemma}
\begin{proof}
For every dependency $D$ with $p(D) > p(u)$, membership in $\Sc{1}(D)$ depends on queue, stage, call, and $p < p(D)$ and
$p < c_D$, all of which $e$ shares with $u$ or satisfies a fortiori. Membership in $\ASc{1}{\ell}(D)$ adds the stage,
the flag and $\ell \in R_B$. Whether $D$ is entered locally depends only on the stage, and an $\mathsf{L}$ ending
compares the source's token, which $e$ and $u$ share. Host order is transitive. Without equal tokens the lemma fails:
an $\mathsf{L}$ chain from $u$ can order an access at $u$'s position that no chain from $e$ orders.
\end{proof}

\begin{proof}[Proof of Lemma~\ref{lem:fast}]
Let $\{e,x\}$ race at $\ell$, with $e$ processed first. We show that \pc{Access}$(x)$ reports a pair that contains $x$.
Let $\varsigma$ be the segment of $x$ that holds $\ell$. Say that $v$ \emph{covers} $e$ for $x$ if $v$ is processed
before $x$, $\ell \in R(v)$, $v$ conflicts with $x$, $v$ and $x$ are not exempt by E1--E4, and $\ordl{\ell}(v,x)$
implies $\ordl{\ell}(e,x)$. A \emph{cover sequence} is $e = v_0, v_1, \dots, v_k$ in which each $v_{j+1}$ covers $v_j$
for $x$, each $v_j$ with $j \ge 1$ covers $e$ for $x$, and, whenever $v_{j+1}$ is a transition, $v_j \in
v_{j+1}.\mathit{hid}$.

\emph{Invariant.} At every point between the processing of $e$ and of $x$, there is a cover sequence whose last element
$v_k$ is \emph{live} for $x$. That is, it is in $\mathit{Wr}_\varsigma$, or in $\mathit{Rd}_\varsigma$ (then $x$
writes), or it is the entry of a shadow record $(v_k, d) \in \mathit{Sh}_\varsigma$ whose dominator $d$ is exempt
from $x$ by E1--E3. Right after $e$ is processed, $k = 0$ and $e$ is in $\mathit{Wr}_\varsigma$ or
$\mathit{Rd}_\varsigma$. Since $e$ and $x$ conflict, $x$ writes if $e$ reads. The invariant can be broken only when an
access $u$ removes $v_k$ from $\varsigma$.
\begin{itemize}
\item (D1) $u$ writes, has no atomic mode, and \pc{Ordered}$(v_k,u) = $ \emph{ok}, so $\ordl{\ell}(v_k,u)$ by Theorem~\ref{thm:sound}. If $x$ is
  not exempt from $u$ by E1--E4, append $u$: it is in $\mathit{Wr}_\varsigma$, it conflicts with $x$, and
  $\ordl{\ell}(u,x)$, a memory order because $u$ writes, gives $\ordl{\ell}(v_k,x)$ by Lemma~\ref{lem:trans}, which
  uses the memory relation if $v_k$ writes and the execution relation if it reads. If $u$ is a transition, $v_k$ is
  added to $u.\mathit{hid}$. If $x$ is exempt from $u$ by E1, E2 or E3, then $u$ has a command, scope or instance, and
  \pc{Drop} records $(v_k, u)$ in the shadow list, so $v_k$ stays live for $x$. $x$ cannot be exempt from $u$ by E4,
  because $u$ has no atomic mode: line~9 of \pc{Access} never drops an entry at an atomic access, except by
  supersession (D2).
\item (D2) $u$ supersedes $v_k$: same queue, stage, flag, kind, scope, instance, load/store marker, atomic mode and
  token, neither a transition nor in a call. By Lemma~\ref{lem:subst}, $u$ covers $v_k$ unless $x$ is exempt from $u$ and
  not from $v_k$. Since the tags for E2--E4 are equal, that can only be E1, and then $u$ has a command and \pc{Drop}
  shadows $v_k$.
\item (D3) $u$ reads with the key $\kappa$ of $v_k$, which includes the token, so $v_k$ reads and $x$ writes. $\mathit{Rd}_\varsigma[\kappa]$
  becomes $[u]$ or $[u, o]$, where $o$ is the first entry of the old list whose command differs from $u$'s. The old list
  is ordered newest first. If $x$ is not exempt from $u$, append $u$ (Lemma~\ref{lem:subst}, execution part). Otherwise,
  since the key fixes E2--E4, $x$ is exempt by E1, so $u$ has a command $c = \mathit{cmd}(x) \ne
  \mathit{cmd}(v_k)$, and $o$ exists. Either $o = v_k$, which then stays, or $o$ is newer than $v_k$ with the
  same key and a command other than $c$, and we append $o$.
\item A shadow record is pruned when its dominator's command, scope and instance are no longer current on its queue.
  By W2, no later access can then be exempt from the dominator by E1--E3, so the record is not live for $x$.
\end{itemize}
In every case each new element covers $e$, by composing the implications.

\emph{Conclusion.} At \pc{Access}$(x)$, $v_k$ is a candidate. Consider \pc{Ordered}$(v_k,x)$. If it is \emph{ok},
then $\ordl{\ell}(v_k,x)$ by Theorem~\ref{thm:sound}, hence $\ordl{\ell}(e,x)$, contradicting the race. If it is not
\emph{ok} and the pair is not an E5 pair, then $(v_k,x)$ is reported. Otherwise line~4 adds
$v_k.\mathit{hid}$, which contains $v_{k-1}$ and overlaps $x$ at $\ell$, and we repeat with $v_{k-1}$. If $v_{k-1}$
was already compared, its verdict is \emph{ok}, which contradicts the race as before, or reported, or R-LT again. The
descent ends at the latest at $v_0 = e$. There \pc{Ordered}$(e,x)$ is not \emph{ok}, and the pair is not an
E5 pair, which would be exempt. So $(e,x)$ is reported. In all cases $x$ is in a reported pair.
\end{proof}

%% -------------------------------------------------------------------
\subsection{Precision (Proposition~\ref{thm:complete})}
\label{app:precision}

\begin{lemma}[Clock completeness]\label{lem:clockc}
Assume that no signal has a stage mask other than $\top$. If $e \reach_x W$ for a processed wait $W$ whose second
stage set has a stage other than $\mathsf{LT}$, then the clock $c$ joined at $W$ has $c[q(e)] \ge p(e)$, and $W$'s relay
holds an unfiltered entry $(t, \bot) \in W.\rho[q(e)]$ with $t \ge p(e)$.
\end{lemma}
\begin{proof}
Induction on the number of semaphore pairs in the chain. The last pair is $\mathsf{Sig} \chainto W$, with $\mathsf{Sig}$
on $q_s$. Every signal has $S = \top$, so \pc{Relay} returns at its line~2 with unfiltered entries only. If $q(e) =
q_s$, then $c[q_s] = p(\mathsf{Sig}) > p(e)$, and the signal's own entry is $(p(\mathsf{Sig}), \bot)$. Otherwise, by
F2, the chain enters $q_s$ through a wait $W'$ that precedes $\mathsf{Sig}$ there and chains onward, so its second set
has a stage other than $\mathsf{LT}$. By induction its clock has $c'[q(e)] \ge p(e)$ and its relay an unfiltered entry
$t' \ge p(e)$. $\mathit{VC}[q_s] \ge c'$ from then on, and $\mathit{RA}[q_s]$ keeps an unfiltered entry for $q(e)$ of at
least $t'$, because $\sqcup$ keeps the largest one; line~2 of \pc{Relay} merges $\mathit{RA}[q_s]$ into $\mathsf{Sig}$'s
relay.
\end{proof}

\begin{proof}[Proof of Proposition~\ref{thm:complete}]
$\Hspec \subseteq \Hpw$ is Theorem~\ref{thm:sound}. For the converse, let $(e,u)$ be reported. They conflict, are not
exempt, and are not an E5 pair, and \pc{Ordered}$(e,u)$ is not \emph{ok}. Suppose $\ordl{\ell}(e,u)$
for every $\ell \in O$. Fix $\ell_1 \in O$ and a chain $C$ for it. By region uniformity (ii), a barrier whose region
contains $\ell_1$ contains $R(e)$ if $e$ is in its first access scope at $\ell_1$, and contains $R(u) \supseteq O$ if $u$
is in its second. If $C$ contains no semaphore pair, it consists of barriers on $q(u)$ and Lemma~\ref{lem:chain}(b,c)
shows that \pc{Chain}$(e)$ m-reaches (x-reaches) its last barrier, so line~4 or~5 returns \emph{ok}. Otherwise let $W$
be the last wait of $C$. By F2, the rest of $C$ consists of barriers on $q(u)$. By Lemma~\ref{lem:clockc}, which applies by (i), $\mathit{VC}[q(u)][q(e)] \ge p(e)$,
so line~7 does not return \emph{race}, and $W$'s relay holds an unfiltered entry at or after $p(e)$, so line~3 of
\pc{ViaWait} examines $W$'s record, and it returns true at line~4 ($E = \top$), line~5 (the wait chain reaches $B$ by
Lemma~\ref{lem:chain}(c)) or line~6 ($s(u) \in E$, or the wait chain reaches the last barrier, which covers $O$).
In each case \pc{Ordered} returns \emph{ok}, a contradiction. Hence $\ordl{\ell}(e,u)$ fails for some $\ell \in O$, and
$\ordl{\ell}(u,e)$ fails too, since by Lemma~\ref{lem:lin} it would have $u$ processed first. So $(e,u)$ races at $\ell$.
\end{proof}

\paragraph{The split trace $\tau^\sharp$}
For each resource $r$, the regions of the barrier structs on $r$ generate a finite Boolean algebra of element sets.
Replace every event $a$ on $r$ by one event $a|_P$ for each atom $P$ of that algebra that meets $R(a)$, with region $R(a)
\cap P$, at consecutive positions and with $a$'s other attributes and tags. Then $\tau^\sharp$ is region-uniform. No
dependency lies between the parts, so no scope or chain changes, and access-scope membership at $\ell$ is that of $a$.
The parts of one access are pairwise disjoint and never conflict. Hence $a|_P$ and $b|_{P'}$ race at $\ell$ in
$\tau^\sharp$ if and only if $\ell \in P \cap P'$ and $a, b$ race at $\ell$ in $\tau$. W1 and W2 carry over.
By the proposition, the pairwise mode on $\tau^\sharp$ reports exactly the parts of the races of $\tau$ when (i) holds.

%% -------------------------------------------------------------------
\subsection{Periodic executions (Theorem~\ref{thm:window})}
\label{app:window}

Write $\mathit{call}(o)$ for the host call that appended $o$, and order calls by host order.

\begin{lemma}[Locality, W.1]\label{lem:loc}
Under P-MONO, every dependency of a chain from $a$ to $b$ was appended by a call in $[\mathit{call}(a),
\mathit{call}(b)]$. Hence whether two events are ordered, and whether they race, is the same in $\tau^\infty$ and in
any prefix of it that contains both, and it is preserved by any map that preserves queues, program order,
attributes, tags and pairing on the calls of that interval.
\end{lemma}
\begin{proof}
On one queue, operations are appended in host order. By P-MONO, a wait is appended after its signal. So every chain
step goes forward in host order, as do $a \in \Sc{1}(D_1)$ and $b \in \Sc{2}(D_n)$. The event cut of an event wait
$D$ matters only through whether a set lies between $a$ and $D$, which is again inside the interval. Under P-MONO no
wait of a prefix lacks its signal, so the prefix drops no operation.
\end{proof}

\begin{lemma}[Retirement, W.2]\label{lem:ret}
Let $a$ and $b$ be events, neither on a presentation queue, with $\mathit{call}(a)$ before $\mathit{call}(b)$, $b$ in a
frame $j \ge w$ and $\mathit{fr}(a) \le j - N$. Then $a \xspec b$ and $a \mspecl{\ell} b$ for every $\ell$.
\end{lemma}
\begin{proof}
If $a$ is on the host, then $b$ is on the host (host order) or in a batch submitted after $a$. In the latter case, the
host signal of that submission has $a$ in its first scope, and the submission wait has $b$ in its second (transitions
included). If $a$ is on a device queue $q$, then by the definition of $N$ the host has waited, before the first
submission or host access of frame $j$, on a $\top$ signal of $q$ that follows $a$. That signal holds $a$ in its first
scope, it chains into the host wait, and the host wait chains into every later host dependency, in particular into
the submission of $b$'s batch, whose wait holds $b$, or it holds the host event $b$ itself. All links are signals and
waits, whose access scopes are their synchronization scopes.
\end{proof}

\begin{lemma}[Presentation, W.3]\label{lem:pres}
Assume P-WSI and P-MONO. Let $b$ be a device write to a swapchain image $i$, held under the acquire $A_b$ and ended
by the present $P_b$. Let $\rho_A$ be the presentation-engine read appended at $A_b$ and $\rho_P$ the one appended at
$P_b$. If $b$ races an engine read $\pi$ of $i$, then $b$ races $\rho_A$ or $\rho_P$.
\end{lemma}
\begin{proof}
All operations of the presentation queue of $i$ are reads of $i$ and signals and waits of every stage. Between $\rho_A$
and the operations of $P_b$, the queue holds only $A_b$'s signals, because an image is not acquired twice without a
present in between. If $\pi$ precedes $\rho_A$ on that queue or is $\rho_A$: $\rho_A \xspec b$ fails, because otherwise
the chain starts at a signal of every stage after $\rho_A$, which also holds $\pi$, and so $\pi \xspec b$.
$b \mspecl{\ell} \rho_A$ fails by Lemma~\ref{lem:loc}, since $\rho_A$ is appended before $b$. So $b$ races $\rho_A$.
Otherwise $\pi$ is $\rho_P$ or follows it. $b \mspecl{\ell} \rho_P$ fails, because otherwise the chain ends at a wait
of every stage before $\rho_P$, which also holds $\pi$. $\rho_P \xspec b$ fails by Lemma~\ref{lem:loc}. So $b$ races
$\rho_P$.
\end{proof}

\begin{proof}[Proof of Theorem~\ref{thm:window}]
Let $\{a,b\}$ race in $\tau^\infty$, with $\mathit{call}(a)$ first. \emph{Distance (W.4).} If neither is a
presentation-engine event, then $\mathit{fr}(b) - \mathit{fr}(a) \le N - 1$ by Lemma~\ref{lem:ret}, or both lie in
the first $w$ frames. Otherwise, by P-WSI one of them is an engine read and the other a device write $b'$ to the same
image. By Lemma~\ref{lem:pres}, $b'$ races $\rho_A$ or $\rho_P$, which lie within $G$ frames of $b'$ by the definition
of $G$. In both cases there is a witness race $\{a',b'\}$ of frame distance at most $D - 1$, which has an event in
common with $\{a,b\}$. \emph{Translation.} While $\mathit{fr}(a') \ge w + p$, the calls from $\mathit{call}(a')$ to
$\mathit{call}(b')$ lie in frames $\ge w + p$, so $\sigma^{-1}$ maps them onto frames $\ge w$, preserving everything
Lemma~\ref{lem:loc} needs, and $\{\sigma^{-1}a', \sigma^{-1}b'\}$ is a race. After finitely many steps, $\mathit{fr}(a')
\le w + p - 1$ and $\mathit{fr}(b') \le w + p + D - 2 = M - 1$. By Lemma~\ref{lem:loc}, this witness is a race of the
first $M$ frames, which are those of $\tau$. Conversely, a race of the first $M$ frames is a race of $\tau^\infty$ by
Lemma~\ref{lem:loc}. The last sentence of the theorem follows with Theorem~\ref{thm:sound} and
Corollary~\ref{cor:sched}.
\end{proof}

\paragraph{Why a fence does not bound the presentation engine}
In the encoding, the engine's read of a presented image is the pair $\rho_P$ (after the present's waits) and $\rho_A$
(at the next acquire of that image, after it waits on the present's completion). A frame fence is a signal on a device
queue. Neither read is in its first scope, because no chain runs from the presentation queue to the device except
through an acquire's signal. So a fence wait proves nothing about them, and $N$ cannot bound them. This is the
reason for the separate parameter $G$ and for Lemma~\ref{lem:pres}.

%% -------------------------------------------------------------------
\subsection{Elision (Theorem~\ref{thm:elide})}
\label{app:elide}

$F^*$ splits each transition $t_B$ into a struct $B_1$ with $B$'s first stage, access and region masks and a second
stage set $\{\mathsf{LT}\}$, a transition use $\lambda_B$ at stage $\mathsf{LT}$, and a struct $B_2$ with first stage
set $\{\mathsf{LT}\}$ and $B$'s second masks. All three have $B$'s call and struct identity. $F^*$ also turns load and
store operations into accesses without the E3 exemption. Identify $\lambda_B$ with $t_B$.

\begin{lemma}[Conservativity of $F^*$]\label{lem:fstar}
For events $a,b$ of $\tau$ and every $\ell$, $\ordl{\ell}(a,b)$ in $F^*(\tau)$ implies $\ordl{\ell}(a,b)$ in $\tau$.
Hence $\Hspec(\tau) \subseteq \Hspec(F^*(\tau))$ under the identification.
\end{lemma}
\begin{proof}
$\mathsf{LT}$ lies in a stage set only if the set names it, so $\top$ does not contain it. $B_1$'s second set
therefore meets only $B_2$'s first set, and conversely. $B_1 \chainto B_2$ is allowed, because they share their
struct identity. $D \chainto B_1$ if and only if $D \chainto B$, and $B_2 \chainto D$ if and only if $B \chainto D$.
For events other than transitions, $\Sc{1}(B_1) = \Sc{1}(B)$, $\Sc{2}(B_2) = \Sc{2}(B)$, and the same holds for the
access scopes. Replacing each $B_1 B_2$ in a chain of $F^*(\tau)$ by $B$ gives a chain of $\tau$ with the same end
scopes. $\lambda_B$ is in the scopes of $B_1$ (second) and $B_2$ (first) only, and $t_B$ is in both scopes of $B$, so
orderings of $\lambda_B$ map to orderings of $t_B$. For the races: $F^*$ removes exemptions and adds events, but adds
no ordering among the events of $\tau$.
\end{proof}

\begin{proof}[Proof of Theorem~\ref{thm:elide}]
Let $\{a,b\} \in \Hspec(\tau_P)$. The plan removes or narrows only barrier structs and keeps every transition, so the
events of $\tau_P$ are those of $\tau$, identified by (call, command, access index). By Lemma~\ref{lem:fstar},
$\{a,b\} \in \Hspec(F^*(\tau_P))$. By Theorem~\ref{thm:sound} applied to $F^*(\tau_P)$, $\{a,b\} \in
\Hpw(F^*(\tau_P))$. By the premise, $\{a,b\} \in \Hpw(F^*(\tau))$. The premise is checked on the final plan, so the
order of the greedy search is irrelevant.
\end{proof}


%% -------------------------------------------------------------------
\subsection{Repair monotonicity (Lemma~\ref{lem:repair})}
\label{app:repair}

\begin{proof}[Proof of Lemma~\ref{lem:repair}]
\emph{Events.} Widening masks adds no access, and an inserted barrier without a transition adds a dependency but no
access event; layouts and regions are unchanged, so the transitions are the same. Hence $\tau$ and $\tau_r$ have the
same access events, identified by (call, command, access index), with the same resources, regions and kinds.

\emph{Scopes grow.} Write $M \sqsubseteq M'$ when every mask of a dependency is contained in the corresponding mask of
its repaired version. The execution scopes $\ESs$, $\ESd$ and the access scopes $\AS$ are monotone in the masks
(each is a union over named stages or flags, possibly expanded to logically earlier or later stages). So for the
repaired dependency $D'$ of $D$: $\Sc{1}(D) \subseteq \Sc{1}(D')$, $\Sc{2}(D) \subseteq \Sc{2}(D')$, and the same for
the access scopes. All other dependencies are unchanged.

\emph{Chains persist.} $D_1 \chainto D_2$ requires that the second synchronization scope of $D_1$ meets the first of
$D_2$ (and, within one call, the struct-identity condition, which a repair does not change). Both scopes only grow,
so every chain of $\tau$ is a chain of $\tau_r$, with end scopes that contain the original ones. An inserted barrier
only adds dependencies and therefore chains.

\emph{Orderings persist.} Execution ordering $\xspec$ and memory ordering $\mspec$ of two events are defined by the
existence of a chain whose first scopes contain the earlier event (and, for $\mspec$, whose first access scope
contains it) and whose last scopes contain the later one. By the two previous steps, every witness in $\tau$ is a
witness in $\tau_r$. Queue-order, semaphore, fence and host orderings are untouched. So $a \spec b$ in $\tau$ implies
$a \spec b$ in $\tau_r$.

\emph{Races.} A race is a conflicting pair, at least one a write, with overlapping regions, not ordered either way.
Conflict and overlap depend only on the events, which are the same, and every ordered pair of $\tau$ stays ordered.
Hence every race of $\tau_r$ is a race of $\tau$: $\Hspec(\tau_r) \subseteq \Hspec(\tau)$.

\emph{Implicit dependencies.} A render pass without an explicit $\mathsf{EXTERNAL}$ dependency has an implicit one.
The grouper widens it to a dependency whose masks contain the implicit masks, which is a widening in the sense above.
A repair that replaced the implicit dependency by an explicit one with fewer stages or flags would not be; the
grouper never emits one.
\end{proof}

%% -------------------------------------------------------------------
\subsection{Minimal frames}
\label{app:counter}

Each frame below is a regression test of the released checker. Each isolates one case in which a plausible
shortcut in the algorithm would accept a race; the frames show why the corresponding lines of
Figure~\ref{fig:alg} are needed. For each we give the pairwise result, the epoch result and the result of the
executable reference semantics.
\begin{itemize}
\item \textbf{Relay through a masked signal.} Compute: write \code{buf} (compute shader), $\mathsf{Sig}(a)$. Graphics:
  $\mathsf{Wait}(a, \text{vertex shader})$, $\mathsf{Sig}(b, \text{compute shader})$. Transfer: $\mathsf{Wait}(b)$,
  read \code{buf} (copy). The wait's second set does not meet the signal's first set, so no chain exists and the
  pair races. Relaying the compute queue's clock entry unfiltered through $\mathsf{Sig}(b)$ would accept it. The
  checker builds the relay of $\mathsf{Sig}(b)$ only from waits whose chain meets its first scope (\pc{Relay},
  line~4), so all three report the race. The host relay (a masked signal waited on by the host, which then submits)
  is handled in the same way.
\item \textbf{Host-signalled wait.} Host: submit, write \code{buf}, then signal value $t{=}5$ from the host. Graphics:
  submission wait, $\mathsf{Wait}(t{=}5, \text{fragment shader})$, read \code{buf} (vertex shader, uniform read).
  The write is ordered only into the fragment stage, so the pair races. Treating the wait as a submission would
  accept it; \pc{ViaWait} has no such shortcut (line~1 uses only full waits). All three report the race.
\item \textbf{Atomic stand-in.} One queue: fill \code{counter}; barrier transfer $\to$ compute; compute atomic
  add; fragment atomic add. The fill and the fragment add race. If the epoch mode dropped the fill at the compute
  add, the fragment add, which is exempt from the compute add by E4, would be compared with nothing. The checker does
  not drop an entry at an atomic access (\pc{Access}, line~9), so all three report the race.
\item \textbf{Token stand-in.} One render pass: a fragment-shader sampled read of image $A$, a fragment-shader
  input-attachment read of $A$ (position token $\pi$), a by-region barrier from the input-attachment read to the
  color-attachment write, and a color-attachment write of $A$ (token $\pi$). The local link orders the input read
  before the write, but not the sampled read, which has no token, so the sampled read and the write race. If the two
  reads shared an epoch read key, the input read would stand in for the sampled read and the race would be missed.
  The read key and \pc{Supersedes} include the token (\pc{Access}, line~10), so both modes and the reference report
  the race, whether or not the two reads come from one command.
\item \textbf{Host write after a fence only.} Device: write \code{buf} (transfer). Host: wait on the fence, then
  write \code{buf} through a mapped pointer. Without a barrier whose second scope holds a host access, the host write
  relies on premise A-HOST (Definition~\ref{def:hostdom}); both modes emit the warning
  \code{timing.premise-host-domain}, and the frame is not guarantee-eligible. Adding a barrier from the transfer
  stage to the host stage with \code{HOST\_WRITE} removes the warning.
\item \textbf{Region cut.} Write $[0,100)$; global barrier; writes $[0,50)$ and $[50,100)$; one call with two buffer
  structs for $[0,50)$ and $[50,100)$ to the fragment shader; read $[0,100)$. No race. Pairwise reports the first
  write and the read, because no struct covers $[0,100)$. The epoch mode and the reference report nothing. This is the
  precision limit of Proposition~\ref{thm:complete}, which errs in the report direction.
\end{itemize}
